\chapter{\B{Laboratory files and terms}{Archivos de laboratorio y términos}}
\label{app:labs}

\section{\B{Reproduce the classroom calculations}{Reproduce los cálculos de aula}}

\B{You can download the files below from the link at the front of the book. Work through the table examples with paper first. You can also use the Python files after you download them. CSV files are tables you can open in a spreadsheet. Their first row gives the column names. Read dates as year, month and day. The ferry notices are for practice.}
{Puedes descargar los archivos del enlace al principio del libro. Empieza por hacer en papel los ejemplos de los cuadros. También puedes usar los archivos de Python después de descargarlos. Los archivos CSV son tablas que puedes abrir en una hoja de cálculo. La primera fila da los nombres de las columnas. Lee las fechas en el orden año, mes y día. Los avisos de ferri son de práctica.}

\noindent\begin{tabularx}{\textwidth}{@{}L{.41\textwidth}Y@{}}
\toprule
\B{File}{Archivo} & \B{Contents}{Contenido}\\
\midrule
\texttt{data/patty-sites.csv} & \B{Fictional site IDs, lunch counts, JMD rents and hub distances.}{ID, conteos, alquileres en JMD y distancias ficticios.}\\
\texttt{data/language-samples.json} & \B{Teaching utterances and two model-free segmentations.}{Enunciados didácticos y dos segmentaciones sin modelo.}\\
\texttt{data/ferry-notices.csv} & \B{Fictional notice records for the retrieval lab.}{Avisos ficticios del laboratorio de recuperación.}\\
\texttt{data/shipment.csv} & \B{Fictional capacities and weights for the dispatch exercise.}{Capacidades y pesos ficticios para el despacho.}\\
\texttt{data/regional-observations.csv} & \B{CTO and STATIN displayed observations. See Chapter~\ref{ch:capstone} for units and source limits.}{Observaciones mostradas por la CTO y STATIN. Véanse unidades y límites en el capítulo~\ref{ch:capstone}.}\\
\texttt{data/place-lab-cases.json} & \B{Fictional lookup fixtures and answer IDs for the geographic index.}{Registros ficticios de consulta e identificadores de respuesta del índice geográfico.}\\
\texttt{checks/verify\_examples.py} & \B{Recomputes selected results, tests the lookup and checks deliberate failures.}{Recalcula resultados, prueba la consulta y comprueba fallos deliberados.}\\
\bottomrule
\end{tabularx}

\B{From the companion directory, use \texttt{python3 checks/verify\_examples.py}. Change the weights in a separate copy of the file, print the scores and record which constraint or priority changed. For the ferry case, use the provided lookup tests before adding a language model. Keep a copy of the original files so you can reproduce the classroom answer. The companion review record uses BootLoops procedures for independent numerical checks, deliberate failures and reference checking \citep{boot}.}
{Desde el directorio complementario, usa \texttt{python3 checks/verify\_examples.py}. Cambia los pesos en una copia, imprime las puntuaciones y anota qué restricción o prioridad cambió. Para el ferri, usa las pruebas de consulta antes de añadir un modelo de lenguaje. Conserva una copia de los archivos originales para reproducir la respuesta de aula. El registro de revisión usa procedimientos de BootLoops para cálculos independientes, fallos deliberados y comprobación de referencias \citep{boot}.}

\section{\B{Keep an evidence record}{Conserva un registro de evidencia}}

\B{Keep a note of where you got each fact. Write the source title, the office that sent it and the place it covers. Write when it came out, when it began to apply, which version you used and what unit it uses. If you work out a new number, save the old numbers and your steps. If a notice changes, keep both copies and mark which one took over. Check that an older question still finds the notice that was in use then.}
{Anota de dónde salió cada dato. Escribe el título de la fuente, la oficina que la envió y el lugar que cubre. Indica cuándo salió, cuándo empezó a valer, qué versión usaste y en qué unidad está. Si calculas una cifra nueva, guarda las cifras originales y tus pasos. Si cambia un aviso, conserva las dos copias y marca cuál reemplazó a cuál. Comprueba que una pregunta sobre una fecha anterior todavía encuentre el aviso que regía en ese momento.}

\section{\B{Terminology used in the two editions}{Terminología de las dos ediciones}}

\small
\begin{longtable}{@{}L{.27\textwidth}L{.26\textwidth}L{.39\textwidth}@{}}
\toprule
English & Español & \B{Use}{Uso}\\
\midrule\endhead
Attention & Atención & \B{Weight and mix information across positions.}{Ponderar y mezclar información entre posiciones.}\\
Context window & Ventana de contexto & \B{Positions the model interface can process in a request.}{Posiciones que admite la interfaz por petición.}\\
Decoder & Decodificador & \B{Module that predicts target tokens using its allowed context.}{Módulo que predice tokens objetivo con el contexto permitido.}\\
Embedding & Representación vectorial & \B{Learned vector for a token, passage or other object.}{Vector aprendido para un token, pasaje u objeto.}\\
Fine-tuning & Ajuste fino & \B{Update model parameters using a specified objective and data.}{Actualizar parámetros con un objetivo y datos definidos.}\\
Grounding & Respaldo en fuentes & \B{Tie an answer's factual claims to applicable evidence.}{Vincular afirmaciones factuales a evidencia aplicable.}\\
Inference & Inferencia & \B{Apply a trained model to a new input.}{Aplicar un modelo entrenado a una entrada nueva.}\\
Logit & Logit & \B{Score before converting a vector to probabilities.}{Puntuación previa a convertir un vector en probabilidades.}\\
Pre-training & Preentrenamiento & \B{Fit shared parameters before task adaptation.}{Ajustar parámetros compartidos antes de adaptar a tareas.}\\
Prompt & Instrucción o entrada & \B{Text or other input given to a model for a task.}{Texto u otra entrada dada al modelo para una tarea.}\\
Retrieval & Recuperación & \B{Find applicable records or passages for a request.}{Encontrar registros o pasajes aplicables a una petición.}\\
Token & Token & \B{Vocabulary item used by a model to represent input or output.}{Elemento del vocabulario con que el modelo representa entradas o salidas.}\\
\bottomrule
\end{longtable}
\normalsize
